% ==============================================================================
% SECTION 6 — CONCLUSION
% ==============================================================================
\section{Conclusion}
\label{sec:conclusion}

We extended the \citet{Willis2025_llm_ipd} \LLM evolutionary
game-theory benchmark to four next-generation frontier models spanning three
providers---48 Moran process conditions, 4,800 simulation runs, identical
experimental protocol.

The headline result is that cooperative bias survives the transition to a new
model generation: nine of twelve model--prompt combinations favour cooperative
equilibria in balanced noiseless conditions, and within the Anthropic lineage
Claude~4.6 is virtually indistinguishable from its predecessor.
What differentiates outcomes is not which generation a model belongs to but
which provider trained it. Gemini~2.5~Flash reaches up to 70\% aggressive
equilibria; GPT-5.4~Mini reaches up to 68\% cooperative under Self-Refine.
That cross-provider gap is larger than any intra-lineage shift we observe,
pointing to alignment objectives---not raw capability---as the dominant
determinant of equilibrium behaviour.

The Self-Refine prompt sharpens this picture in a troubling direction.
It narrows the aggressive--cooperative payoff gap (\ICD rises toward 1.0)
consistently across all four models, including Claude~4.6 Refine at \ICD = 0.913.
Prompt engineering designed to improve strategy quality may therefore
inadvertently reduce the cooperative advantage, raising the equilibrium
aggression rate in deployed multi-agent populations.
Noise robustness, meanwhile, has not improved: Claude~4.6 shows
$\Dnoise \approx 12$~pp, essentially matching the 13~pp
reported for Claude~3.5~Sonnet in the original benchmark, while biased noisy
conditions converge to aggressive equilibria in the 53--71\% range,
with only 2 of 12 configurations exceeding the 67\% prior.

These findings raise questions the current design cannot answer.
Do the cross-provider differences reflect divergent RLHF reward signals,
differences in training corpora, or systematic variation in post-training
alignment objectives? Do cooperative equilibria persist in $n$-player or
asymmetric games, or are they an artefact of the symmetric two-player IPD
structure? And would mixed-provider populations---where a Gemini agent
competes against a GPT agent---generate equilibrium dynamics that neither
provider's models produce in isolation? The persistence of noise sensitivity
across two Anthropic generations suggests that robustness may require
targeted interventions beyond general capability improvements.

Having this baseline makes those questions tractable.
The benchmark protocol and complete strategy datasets are reproducible
and extensible; % VERIFY: confirm repository URL before submission
we release simulation code and strategy libraries so that each new
frontier model can be slotted into the same 48-condition design,
turning this cross-provider snapshot into a running longitudinal record.

\section*{Acknowledgements}
[Acknowledgements to be added before submission.]
